% Capítulo 1

\chapter{Estrutura do traballo}

\label{Chapter1} % para referenciar este capítulo en calquera parte deste documento, emprega \ref{Chapter1} 



\section{Preliminares}
Recoméndase consultar a norma UNE 50136:1997: documentación, presentación de tese e documentos similares. Tamén existen recursos en liña que axudan a organizar todo o proceso de elaboración do Traballo. 

Ten en conta, que a morfoloxía do documento non está condicionada pola plantilla e podes estruturalo acorde ao tipo de documento que vaias a realizar (proxecto construtivo, traballo técnico ou de investigación).

\subsection{Cuberta}
Pode utilizarse como plantilla a cuberta que se inclúe neste asegurándose de que se inclúen os datos precisos. Debe conter todos os datos que se especifican: nome completo do autor, título do traballo, nome completo do titor, data, título ao que se opta, etc.
\subsection{Portada}
A cuberta é a tapa do documento e é a primeira folla que aparece ao abrir a tapa do documento. Igualmente pode utilizarse como modelo a que se inclúe aquí e debe conter todos os datos que figuran nesta: o nome da Escola, a titulación e, no seu caso, a intensificación, o título do TFG/TFM, os nomes do autor, do titor e, no seu caso, do relator, a área de coñecemento, o departamento e o ano de execución do traballo.

\subsection{Resumo}
Texto que mostra de forma abreviada e precisa o contido do traballo. Deberase incorporar unha versión en galego, castelán e inglés que se colocarán en orde consecutivo. 
\subsection{Prefacio ou Introdución}
Debe incluír unha breve explicación das razóns que levaron á realización do Traballo, o propósito e os obxectivos que se pretenden, o ámbito, alcance e límites da investigación, así como a metodoloxía empregada e, se se considera oportuno, un avance das conclusións alcanzadas. 
\subsection{Índice xeral}
Táboa de contidos onde se reflicten todas as partes do Traballo e os seus anexos, se os houbese. Deben aparecer os títulos, na súa orde e con indicación da páxina na que se poden atopar.
Se o Traballo consta de varios volumes, cada un deberá levar o seu propio índice, pero débese incluír tamén un índice xeral.
\subsection{Lista de ilustracións e táboas}
Se o Traballo inclúe ilustracións e táboas pódese engadir unha listaxe que inclúa o número identificativo que figura dentro do texto, a lenda e o número da páxina na que se atopa.
Tamén é conveniente mencionar os datos sobre as fontes de onde se obtiveron ditas ilustracións se non se incluíron no propio texto da memoria.
\subsection{Lista de abreviaturas e símbolos}
O Traballo debe conter as abreviaturas e símbolos internacionalmente recoñecidos. Se se incorporan unidades, abreviaturas ou acrónimos que poidan ser pouco coñecidos deberanse explicar brevemente nestas listas.
\subsection{Glosario}
Os termos que requiran definición ou explicación deberanse incorporar nun glosario.

\section{Texto principal}
Para a súa elaboración é importante ter en conta a súa distribución en capítulos e seccións numeradas. 
O texto debe comezar cunha introdución que mostre as investigacións previas existentes sobre o tema e destacar os obxectivos e métodos seguidos para levar a cabo a investigación ou análise do tema tratado.
Para finalizar débense escribir as conclusións que deben estar en relación cos obxectivos marcados previamente.

\section{Bibliografía}
Debe conter as referencias bibliográficas dos documentos consultados para demostrar as bases do traballo realizado e avalar os datos incorporados e citados no texto.
Elaborarase de forma normalizada, para o que se aconsella utilizar a norma UNE vixente (actualmente a UNE ISO 690:2013. Información e documentación. Directrices para a redacción de referencias bibliográficas e de citas de recursos de información).
Para a elaboración desta parte do Traballo recoméndase consultar a Web, guías e axudas para a elaboración das referencias bibliográficas.

\section{Anexos}
Pódese incluír desta forma material extenso utilizado no traballo, importante para xustificar os resultados e as conclusións obtidas, pero que non é esencial para a comprensión do texto principal. Poden ser datos estatísticos, lexislación, etc.
A paxinación debe ser correlativa e continuar a do texto principal. Cada un dos anexos debe identificarse cunha letra maiúscula do alfabeto, comezando pola letra A, precedida da palabra Anexo.
